problem:  
Let \((X,d,\mathfrak{m})\) be a complete, non-collapsed \(RCD(K,N)\) space with \(K\in\mathbb{R}\), \(1<N<\infty\).  
Assume that there exists an integer \(k_0\) such that, for every \(x\in X\), every measured tangent cone \(T_xX\) (in the pointed measured Gromov–Hausdorff sense) is uniquely Euclidean, i.e.  
\[
T_xX \cong (\mathbb{R}^{k_0},|\cdot|,\lambda^{k_0})
\]
as a metric measure space.  

**Question:**  
Must \(X\) then be globally isomorphic (as a metric measure space) to a \(C^{1,\alpha}\) Riemannian manifold \((M^{k_0},g,\mathrm{vol}_g)\) with \(\mathrm{Ric}_g\ge K\)?  
Equivalently, can a complete, non-collapsed \(RCD(K,N)\) space admit *Euclidean tangent cones at every point* yet fail to be a smooth Riemannian manifold?  

Why is it a "good" problem:  
This problem isolates the deepest unresolved boundary between **synthetic Ricci curvature geometry** and **classical smooth Riemannian geometry**. Current theory (Cheeger–Colding, De Philippis–Gigli, Mondino–Naber) ensures that non-collapsed \(RCD(K,N)\) spaces are rectifiable with Euclidean tangents almost everywhere, but singularities of codimension ≥ 4 may remain. The question probes whether the *universal* Euclideanity of tangents removes these singularities entirely—whether the local Riemannian structure implied by the tangents extends to a global \(C^{1,\alpha}\) manifold structure.

A positive resolution would establish that synthetic Ricci bounds and infinitesimal Hilbertianity fully determine the smooth category, completing the bridge from \(RCD\) spaces to Riemannian manifolds. A counterexample would reveal genuinely nonsmooth phenomena invisible to curvature-dimension conditions, reshaping our understanding of the fine structure of \(RCD\) geometry. The problem is stated with striking simplicity but touches profound issues of regularity, Gromov–Hausdorff stability, and analytic versus geometric reconstruction—capturing exactly the “narrow entrance, profound depth” ideal of a good mathematical problem.